% =====================================================================
%  Cloud removal from satellite images - plain-language guide
%  Companion to ecrformer_pipeline_technical.tex
%  Compile with: pdflatex ecrformer_pipeline.tex   (run twice)
% =====================================================================
\documentclass[11pt,a4paper]{article}

\usepackage[margin=1in]{geometry}
\usepackage{amsmath}
\usepackage{array}
\usepackage{float}
\usepackage{xcolor}
\usepackage{listings}
\usepackage[hidelinks]{hyperref}

\lstset{
  basicstyle=\ttfamily\footnotesize,
  breaklines=true,
  columns=fullflexible,
  frame=single,
  framerule=0.3pt,
  rulecolor=\color{gray!50},
  backgroundcolor=\color{gray!8},
  showstringspaces=false,
  language=bash
}

\newcommand{\code}[1]{\texttt{#1}}
\newcommand{\term}[1]{\textbf{#1}}

\title{\textbf{Removing Clouds from Satellite Photos}\\[0.3em]
\large A Plain-Language Guide:\\What We Built, What We Measured, and What Is Going Wrong\\[0.2em]
\normalsize Written 2026-10-15}
\author{Vaishnav (yovyshh)}
\date{2026-10-15}

\begin{document}
\maketitle

\begin{abstract}
\noindent
Satellite photos of the ground are often ruined by clouds. We built a computer program that tries
to paint over the clouds and show what was underneath. This document explains --- in everyday
language, without assuming you know any machine learning --- what the program is, what data it
learns from, how well it is doing, and the real problems we have found with it. Every number in
here came from a script we actually ran. Where a number is unknown, we say so instead of guessing.

A technical version of the same material, with the full mathematics, is in
\code{ecrformer\_pipeline\_technical.tex}.
\end{abstract}

\tableofcontents

% =====================================================================
\section{The problem, in one paragraph}
% =====================================================================

Imagine you photograph a field. A big cloud drifts over it and your photo is mostly grey fluff.
You cannot wait for another day, because the crops change, the season changes, and the cloud may
come back.

But you have a second camera that is not bothered by clouds at all: a \term{radar}. Radar sends out
its own radio waves and listens for the echo, the way a bat does. It works at night, in fog, and
right through cloud. The catch is that radar pictures are grey, grainy and hard for a human to read
--- no colour, no nice texture.

So the plan is: feed a computer (a) the cloudy colour photo, (b) the radar picture of the same
spot, and (c) a \term{stencil} that marks which pixels are cloudy. Ask it to paint the missing
parts. That is the whole job. The program we are working with is called \term{ECRformer}.

% =====================================================================
\section{The data we learn from}
% =====================================================================

We do not have perfect answers to copy, so we use pairs of images of the same piece of land taken
on different days: one cloudy, one clear. The clear one is the ``answer sheet'' the computer is
marked against.

\begin{center}
\begin{tabular}{@{}lp{9.5cm}@{}}
\hline
\textbf{Thing} & \textbf{What it means} \\
\hline
Training pairs & 33{,}468 image pairs used for learning \\
Validation pairs & 1{,}115 pairs used for testing during development \\
Final exam pairs & a whole scene kept locked away and never touched \\
Patch size & $256 \times 256$ pixels (a small square of land) \\
Colour channels & 4: red, green, blue, and near-infrared \\
Radar channels & 2: called VV and VH (two ways of listening to the echo) \\
Cloud stencil & 1 more channel, marking cloud \\
\hline
\end{tabular}
\end{center}

\subsection{What ``one image'' actually looks like inside the computer}

This is where matrices come in. A single colour patch is \textbf{not} one picture --- it is
\textbf{four grids of numbers}, one per channel, each grid $256 \times 256$. Every number is a
brightness between 0 and 1 (0 = black, 1 = brightest). The fourth grid, near-infrared, is invisible
to your eye but very useful for telling plants from bare ground.

\begin{center}
\begin{tabular}{@{}lll@{}}
\hline
Grid & Size & Meaning \\
\hline
Red & $256 \times 256$ & how much red light bounced back, per pixel \\
Green & $256 \times 256$ & how much green light \\
Blue & $256 \times 256$ & how much blue light \\
Near-infrared & $256 \times 256$ & invisible to us, but good for plants \\
Radar VV & $256 \times 256$ & echo strength, one listening mode \\
Radar VH & $256 \times 256$ & echo strength, other listening mode \\
Cloud stencil & $256 \times 256$ & 0 = clear sky here, 1 = cloud here \\
\hline
\end{tabular}
\end{center}

So the computer receives seven grids at once and must produce four grids back (the repaired red,
green, blue and near-infrared). Keeping the four grids \emph{pixel-aligned} matters: the 500th
pixel of the red grid must describe the same 500th spot on the ground as the 500th pixel of the
radar grid. The data was prepared that way in advance, so we do not have to fix it.

% =====================================================================
\section{How we prepare each image before the model sees it}
% =====================================================================

Three simple jobs, done on the fly every time an image is loaded.

\paragraph{1. Putting the radar on a friendly scale.}
Raw radar echo values are stored in a unit called decibels and run from about $-37$ up to $0$. The
program does not like negative numbers much, so we squeeze them into the range 0 to 1 with a simple
rule: anything below the weakest useful echo counts as 0, anything at the strongest counts as 1, and
everything else is spread evenly in between. This is just a change of units, like converting
Celsius to Fahrenheit --- no information is created or destroyed.

\paragraph{2. Tidying up the cloud stencil.}
The raw stencil is noisy: scattered single pixels claim to be cloud. We clean it in three steps.

\begin{center}
\begin{tabular}{@{}llp{6.6cm}@{}}
\hline
Step & Action & Why \\
\hline
1 & average each $9\times9$ neighbourhood and round it to clear or cloud & removes speckle --- one odd pixel no longer claims to be a cloud \\
2 & grow the cloud region outward by 2 pixels & clouds have soft edges; this catches the fuzzy fringe \\
3 & blur the edges & so the computer gets a soft transition instead of a hard cut-out \\
\hline
\end{tabular}
\end{center}

\paragraph{3. Choosing how much to look at.}
During learning we cut a $128 \times 128$ window out of each patch (always at a random spot, so the
model sees variety). When we test, we use the centre window so every run sees exactly the same
thing. For the final picture quality checks we use the whole $256 \times 256$ patch with nothing cut
off. Section~\ref{sec:ruler} explains why mixing those up caused us trouble.

% =====================================================================
\section{How the model works, in plain words}
% =====================================================================

\subsection{The shape of the machine}

Think of it as a factory with three stages.

\begin{enumerate}
\item \textbf{Two front doors.} One door takes the radar picture; the other takes the colour photo
      plus the cloud stencil. Radar and colour are very different kinds of evidence, so they are not
      mashed together immediately --- each gets its own entrance.
\item \textbf{A committee of pattern-finders.} Behind the doors sits a stack of four stages that
      look for patterns, from small (an edge of a roof) to large (a whole field). The mechanism is
      called a \term{transformer}; practically, it means the machine compares every part of the
      image with every other part and decides what belongs together. Our version is the
      ``light'' one --- deliberately small so it can run on one modest graphics card.
\item \textbf{One painter.} The last stage paints four grids of numbers: the repaired red, green,
      blue and near-infrared. Alongside the main painter there are six small helper painters, three
      belonging to the middle of the factory and three to the end. They are not the final answer;
      they exist so that mistakes can be corrected early rather than only at the very end, the same
      way a maths teacher marks your working, not just your final answer.
\end{enumerate}

\subsection{Why the painter starts with a grey-green blank canvas}

Before training begins, the painter is initialised so that it outputs the \emph{average} colour of
clear land: roughly red $0.085$, green $0.092$, blue $0.113$, near-infrared $0.188$. Asked to paint
nothing in particular, it paints a flat average landscape.

\begin{center}
\begin{tabular}{@{}lllll@{}}
\hline
Band & Red & Green & Blue & Near-infrared \\
\hline
Starting value & 0.0847 & 0.0919 & 0.1127 & 0.1884 \\
\hline
\end{tabular}
\end{center}

This is a sensible trick --- it means the machine starts from a plausible guess instead of random
noise, so it learns faster. But it also plants the seed of our biggest problem
(Section~\ref{sec:wrong}). A painter trained to start from ``the average of everything'' tends to
return to averages when it is unsure.

\subsection{Two different products come out}

This distinction matters more than any other single thing in this document.

\begin{center}
\begin{tabular}{@{}lp{9.4cm}@{}}
\hline
\textbf{Name} & \textbf{What it is} \\
\hline
\term{RAW} (\code{pred\_direct}) & the painter's picture, used as-is, everywhere \\
\term{BLEND} & the painter's picture used \emph{only inside the cloud stencil}; outside the stencil, the original cloudy photo is kept untouched \\
\hline
\end{tabular}
\end{center}

The blend rule in words: \emph{``where there is no cloud, show the original photo unchanged; where
there is cloud, show the painter's picture.''} In symbols, with $M$ the stencil (0 clear, 1 cloud),
$S2_c$ the cloudy photo and $P$ the painter's picture:

\begin{equation}
\hat Y = (1-M)\,S2_c \;+\; M\,P
\end{equation}

Read it as: when $M=1$ (cloud) the left term vanishes and you get $P$; when $M=0$ (clear) the right
term vanishes and you get $S2_c$.

The blend sounds safer, and in one narrow way it is: it \emph{guarantees} clear ground is never
disturbed. But it is not a claim about accuracy --- if the stencil is wrong, the blend faithfully
reproduces the clouds you were trying to remove. We measured this happening
(Section~\ref{sec:wrong}).

% =====================================================================
\section{How we score the results}
% =====================================================================

Four scores matter. Here they are without jargon.

\begin{center}
\begin{tabular}{@{}lp{4.2cm}p{6.0cm}@{}}
\hline
\textbf{Score} & \textbf{Range} & \textbf{What it really tells you} \\
\hline
PSNR & decibels, higher is better & Overall closeness to the answer sheet. Rough guide: above 30 is respectable, 20 is poor. \textbf{Warning:} it is an average, so painting a smooth blurry smear can score well. \\
SSIM & 0 to 1, 1 is perfect & Do the shapes and patterns match? 0.89 means the layout is broadly right. \\
L1 / MAE & 0 to 1, lower is better & On average, how far off is each pixel's brightness? 0.02 means about 2\% off. \\
\term{Gradient} or detail ratio & 1.0 is perfect & \textbf{The one that catches blur.} See the analogy below. \\
Cloud coverage & 0 to 100\% & How much of this patch the stencil marks as cloudy. \\
\hline
\end{tabular}
\end{center}

\paragraph{The detail ratio, explained by analogy.}
Take a photo of a grassy field and blur it heavily. The averaged colour hardly changes --- the field
is still green, roughly the same shade. What you destroyed is the \emph{texture}: the individual
blades, the flecks of light, the fine edges. PSNR will barely notice. The detail ratio is designed
to notice: it compares how much fine detail (pixel-to-pixel change) exists in the picture versus in
the answer sheet.

\begin{itemize}
\item ratio $= 1.0$: the texture is fully preserved;
\item ratio $= 0.6$: only 60\% of the fine detail survives --- the picture is noticeably soft;
\item ratio $= 0.25$: three-quarters of the texture is gone --- it looks like a watercolour.
\end{itemize}

We also report \term{NDVI}, which detects vegetation health, because a cloud-removal method that
wrecks the vegetation signal is not useful to an agronomist even if its PSNR looks fine.

% =====================================================================
\section{Where the model stands today}
% =====================================================================

All numbers below come from the saved baseline --- the best version of the model we had before this
week's experiment, frozen and backed up.

\subsection{Eight carefully chosen test patches, examined at full size}

\begin{center}
\begin{tabular}{@{}lrrrrr@{}}
\hline
Patch & Cloudy \% & Raw PSNR & Blend PSNR & SSIM & Detail ratio \\
\hline
\code{s2\_119\_p208} & 0.0 & 24.94 & 10.76 & 0.80 & \textbf{0.32} \\
\code{s2\_119\_p236} & 0.9 & 22.79 & 10.63 & 0.74 & \textbf{0.26} \\
\code{s2\_119\_p92}  & 0.0 & 27.87 & 13.61 & 0.87 & 0.48 \\
\code{s2\_40\_p268}  & 5.6 & 28.20 & 28.37 & 0.94 & 0.64 \\
\code{s2\_65\_p328}  & 0.0 & 27.46 & 31.14 & 0.94 & 0.77 \\
\code{s2\_65\_p124}  & 0.8 & 42.85 & 41.39 & 0.98 & 0.85 \\
\code{s2\_119\_p100} & 100.0 & 26.51 & 26.51 & 0.90 & 0.64 \\
\code{s2\_65\_p336}  & 89.3 & 34.43 & 34.25 & 0.95 & 0.90 \\
\hline
\textbf{Average} & & \textbf{29.38} & \textbf{24.58} & \textbf{0.89} & \textbf{0.61} \\
\hline
\end{tabular}
\end{center}

Read the average detail ratio: \textbf{0.61}. The model keeps about sixty per cent of the fine
detail. That is the headline finding of this whole body of work.

\subsection{Twelve patches, sorted by how much detail survived}

\begin{center}
\begin{tabular}{@{}lrr@{}}
\hline
Patch & Cloudy \% & Detail kept \\
\hline
\code{s2\_119\_p264} & 5.4 & 0.68 \\
\code{s2\_134\_p484} & 50.2 & 0.68 \\
\code{s2\_119\_p100} & 100.0 & 0.66 \\
\code{s2\_119\_p208} & 0.0 & 0.27 \\
\code{s2\_119\_p236} & 0.9 & 0.23 \\
\hline
\end{tabular}
\end{center}

The pattern is uncomfortable: the worst detail retention is on the patches with the \emph{least}
cloud. Those are the patches where the model has almost nothing to remove --- and it flattens the
genuinely clear ground anyway.

\subsection{Sanity check on an earlier picture we published}

Five patches from a comparison figure were re-measured independently. All five numbers reproduced to
within 0.01 dB, so that figure can be trusted.

\begin{center}
\begin{tabular}{@{}lrrrr@{}}
\hline
Patch & Cloudy \% & Cloudy photo & Raw & Blend \\
\hline
\code{s2\_65\_p88} & 2.4 & 35.37 & 37.40 & 35.38 \\
\code{s2\_3\_p632} & 12.6 & 23.98 & 32.62 & 29.85 \\
\code{s2\_40\_p476} & 35.0 & 14.45 & 28.42 & 27.33 \\
\code{s2\_40\_p368} & 62.5 & 11.66 & 34.57 & 27.31 \\
\code{s2\_3\_p680} & 87.3 & 10.88 & 29.14 & 28.20 \\
\hline
\end{tabular}
\end{center}

% =====================================================================
\section{What is going wrong}
\label{sec:wrong}
% =====================================================================

\subsection{1. The model paints in soft watercolour}

This is the root problem and everything else is downstream of it. The model keeps about 61\% of the
fine detail, and on the hardest patches as little as 23\%. In plain terms: the clouds disappear, the
colours are roughly right, and the picture is soft and slightly plastic.

Why? Because of what we ask it to do. The training rule says: \emph{be close to the answer sheet on
average}. When several different answers are plausible for a pixel --- a field could be grass, or
the same field a month later --- and you are rewarded for being close on average, the safest answer
is the \emph{middle} of the plausible options. The middle of ``grass'' and ``ploughed field'' is a
featureless brown. Survival by averaging is the same reason a committee-written report reads
blandly: nobody's sharp opinion survives the averaging.

Nothing in the original training rule ever gave the model a point for a crisp edge.

\subsection{2. The headline score hides the problem}

The average PSNR is a respectable 29.38 dB. The most impressive-looking numbers sit exactly where
the picture is flattest. Patch \code{p56} reports a spectacular 44.09 dB, and \code{p208} jumps
from 10.76 dB to 24.94 dB --- yet \code{p208} is the second-worst patch of the whole set for detail
retention (0.27). A blurry watercolour can win this exam. If we report only PSNR, we are reporting
the flattering half of the story.

\subsection{3. The cloud stencil is sometimes plainly wrong}

This is the most surprising finding, and it has real consequences. Two patches tell the story.

\begin{center}
\begin{tabular}{@{}lrr@{}}
\hline
Patch & How bad the cloudy photo is & What the stencil claims \\
\hline
\code{s2\_134\_p56} & 8.40 dB (badly clouded) & 100\% cloud \\
\code{s2\_3\_p680} & 10.88 dB (badly clouded) & 87\% cloud \\
\code{s2\_119\_p208} & \textbf{10.76 dB (badly clouded)} & \textbf{0\% cloud} \\
\code{s2\_119\_p236} & \textbf{10.63 dB (badly clouded)} & \textbf{0.9\% cloud} \\
\code{s2\_119\_p100} & 23.45 dB (mildly affected) & 100\% cloud \\
\hline
\end{tabular}
\end{center}

\code{p208} is exactly as cloudy as \code{p680}, yet the stencil says there is no cloud at all. The
second cloud measurement stored alongside it agrees with the stencil, so this is not a display
glitch --- two independent cloud measures are both contradicted by what the pixels actually look
like.

The practical damage: because the stencil says ``no cloud here'', the blend rule returns the
original cloudy photo on those patches. The model had actually improved them by 12 to 14 dB, and the
blend throws that improvement away. It also means any conclusion we draw about ``the hard patches''
is a conclusion about a possibly broken stencil, not about the model.

\subsection{4. The answer sheet was photographed on a different day}

Our clear reference images are of the same place but \emph{not the same moment} --- they come from
another date. So if the land changed in between (crops harvested, fields flooded, season turned),
the difference between the cloudy photo and the ``answer'' is partly \emph{real change}, not cloud.

This matters because the model can score well on such a patch without removing any cloud at all ---
it just paints the typical look of that kind of terrain. That is a very plausible explanation for
the patches where detail retention is worst (0.23--0.27) while PSNR looks decent. It also means our
current scoring cannot fully separate ``removed the cloud'' from ``guessed the season''.

\subsection{5. We have been using two different rulers}
\label{sec:ruler}

Numbers measured on a $128 \times 128$ centre window, counting only the cloudy pixels, are not
comparable to numbers measured on the whole $256 \times 256$ patch counting every pixel. The two
differ by up to 2.3 dB and sometimes in opposite directions. Mixing them --- or quoting one and
comparing it to the other --- would be a genuine error. We now keep the two families of numbers
strictly apart and say which ruler is in use.

There is a second wrinkle: the saved best model stores a score of 30.1982 that was measured under
the \emph{blended} definition, and the record does not say so. Our current default is the RAW
definition. The same field can look better or worse purely because the definition silently changed.

\subsection{6. We cannot cleanly restart the previous training run}

The saved training state does not include the learning-rate schedule's position. Without it, a
restart would silently use the wrong learning rate and quietly waste days of computation. We also
found that the saved ``best'' model is from epoch 30 while a later file sits at epoch 32 --- so
training continued past the best point without improving, and nothing recorded why it stopped. In
practice the project is pinned at the frozen epoch-30 model, and the sensible way forward is a fresh
run that starts \emph{from} it rather than a continuation of the old one.

\subsection{7. Our fake clouds did not look like real clouds}

One idea was to make training harder by adding artificial clouds to clear images, so the model sees
more practice examples. We generated 96 such examples from real cloudy patches.

\begin{center}
\begin{tabular}{@{}lp{8.4cm}@{}}
\hline
\textbf{Attempt} & \textbf{Why we rejected it} \\
\hline
Copy a real cloud's shape onto a clear scene & The brightness was right but the textures came out too smooth, and in near-infrared especially it did not look like a real cloud at all. \\
\hline
Transfer only the difference between a cloudy and clear scene & Far too extreme: 77 of 96 examples produced impossible pixel values outside the 0--1 range, i.e.\ darker than black or brighter than white. \\
\hline
\end{tabular}
\end{center}

So the artificial clouds do not yet behave like the real problem, and training on them would teach
the model to fix the wrong kind of damage. This experiment is shelved until the fakes pass physical
sanity checks.

% =====================================================================
\section{What we are trying right now: the ``keep the edges'' penalty}
% =====================================================================

Since the core problem is that nothing ever rewarded a crisp edge, we added exactly one new rule to
the training: \emph{punish the model whenever its picture is smoother than the answer sheet}. That
is all. The architecture, the data, the optimiser --- everything else is unchanged, so if anything
improves we will know precisely what caused it.

In words: for every pair of neighbouring pixels, compare how much the brightness changes in the
model's picture versus in the answer sheet, and pay a fine for the difference. The total training
penalty becomes ``the old penalty plus a small multiple of the new one''.

\subsection{How big should the fine be? We measured it instead of guessing}

Too small and the new rule does nothing; too big and it ruins the picture chasing noise. Rather than
pick a number by feel, we ran the existing model over eight batches of training images and measured
how big the two penalties already were. We then chose the multiplier that makes the new rule
contribute about \textbf{5\%} of the total --- big enough to matter, small enough not to take over.

\begin{center}
\begin{tabular}{@{}rrrrrr@{}}
\hline
Batch & Main penalty & Encoder helpers & Decoder helpers & New edge penalty & Share of total \\
\hline
0 & 0.0089 & 0.0851 & 0.0163 & 0.0053 & 3.2\% \\
1 & 0.0145 & 0.1069 & 0.0214 & 0.0112 & 4.6\% \\
2 & 0.0114 & 0.0901 & 0.0172 & 0.0094 & 4.8\% \\
3 & 0.0081 & 0.0561 & 0.0113 & 0.0065 & 4.8\% \\
4 & 0.0123 & 0.1351 & 0.0205 & 0.0122 & 5.2\% \\
5 & 0.0080 & 0.0532 & 0.0117 & 0.0072 & 5.4\% \\
6 & 0.0086 & 0.0595 & 0.0126 & 0.0092 & 6.5\% \\
7 & 0.0113 & 0.0759 & 0.0166 & 0.0097 & 5.2\% \\
\hline
\end{tabular}
\end{center}

The chosen multiplier is \textbf{0.079433536}. Median share came out at exactly 5.0\%, worst case
6.5\% --- comfortably inside the 10\% ceiling we allowed ourselves.

\subsection{How we will judge whether it worked}

Decided in advance, before the run, so we cannot move the goalposts afterwards.

\begin{enumerate}
\item The two problem patches (\code{p208}, \code{p236}) must both keep \emph{more} fine detail.
\item No important patch may lose more than 0.5 dB of overall score.
\item The average score over the eight patches may not drop by more than 0.1 dB.
\item Secondary quality must not degrade: structure, average brightness error, colour balance and
      impossible pixel values all stay within tight limits.
\item No crashes, no absurd values, and no invented detail that is not really there.
\end{enumerate}

If any of these fail, the answer is ``this did not work'' and we keep the evidence and fall back to
the frozen baseline. Even a pass would only mean ``worth investigating further'' --- three epochs is
a short screening run, nowhere near proof.

% =====================================================================
\section{Where we are at this moment}
% =====================================================================

The three-epoch experiment is running now, started 08:53 on 2026-10-15 on a single graphics card.

\begin{center}
\begin{tabular}{@{}llrrrr@{}}
\hline
Epoch & Finished at & Minutes & PSNR (raw) & SSIM & L1 \\
\hline
1 & 09:29 & 36.0 & 30.27 & 0.8874 & 0.0199 \\
2 & 10:04 & 35.0 & 30.25 & 0.8874 & 0.0200 \\
3 & running & --- & --- & --- & --- \\
\hline
\end{tabular}
\end{center}

\begin{itemize}
\item The learning-rate schedule behaved exactly as designed (we can see each epoch's value in a
      log: $0.0001447$, then $0.0001111$, then $0.0000437$ --- a smooth glide down).
\item The new edge penalty landed at \textbf{5.04\%} of the training total during the real run,
      matching the 5\% we aimed for.
\item Epoch 2 scored 0.02 dB \emph{lower} than epoch 1. That is a rounding-level difference and not
      yet meaningful; the decisive test is the eight-patch comparison against the frozen baseline,
      which runs as soon as epoch 3 ends.
\end{itemize}

\paragraph{Important caveat.} The 30.27 dB above is measured with the $128$ centre-window,
cloudy-pixels-only ruler over 1{,}115 validation patches. The 29.38 dB baseline figure is measured
with the whole-$256$-patch, every-pixel ruler over 8 patches. \textbf{They are not comparable.} Any
claim of improvement must come from the same ruler on the same patches.

% =====================================================================
\section{Jargon decoder}
% =====================================================================

\begin{center}
\begin{tabular}{@{}lp{9.6cm}@{}}
\hline
\textbf{Word} & \textbf{One-line meaning} \\
\hline
SAR / radar & The cloud-piercing instrument; grey, grainy, but works in any weather \\
Stencil / mask & A map of which pixels are cloudy \\
Blend & ``Keep the original where it is clear, use the painter's work where it is cloudy'' \\
RAW prediction & The painter's picture with no blending applied \\
PSNR & Similarity score in decibels; higher is better; averages away blur \\
SSIM & Structure similarity, 0 to 1; do the shapes match \\
L1 / MAE & Average brightness error; lower is better \\
Gradient / detail ratio & How much fine texture survived, 1.0 is perfect \\
Epoch & One full pass through all 33{,}468 training pairs \\
Patch & A single $256\times256$ square of land \\
Checkpoint & A saved snapshot of the model's learned settings \\
Auxiliary loss & A helper penalty applied early in the network, not at the final output \\
EMA & A smoothed, averaged version of the model; usually the better one to use \\
Over-smoothing & The blurry effect this whole document is about \\
\hline
\end{tabular}
\end{center}

% =====================================================================
\section{What happens next}
% =====================================================================

\begin{enumerate}
\item Finish epoch 3 and run the eight-patch comparison against the frozen baseline; record
      pass or fail against the five criteria above.
\item Settle the \code{p208}/\code{p236} question: open the original upstream data and check whether
      those two patches were paired and masked correctly. Until we know, any conclusion drawn from
      them is unsafe --- and unfortunately they are exactly the patches criterion 1 depends on.
\item Pick one ruler and stick to it in all reporting, so numbers can never be silently compared
      across protocols.
\item Decide whether fixing the cloud stencil is worth doing before further training, since it
      currently decides whether the blended product is usable at all.
\item Record the scoring definition and the schedule position inside every future checkpoint, so a
      run can be resumed without guesswork.
\end{enumerate}

\vfill
\noindent\rule{\textwidth}{0.3pt}\\
\footnotesize
Every number in this document was produced by a script in the project. Values that are not recorded
anywhere --- the original run's exact penalty settings, its scoring definition, and the schedule
position of the old checkpoints --- are stated as unknown rather than estimated. The locked-away
final-exam scene was never opened at any point described here.

\end{document}
